\section{生态系统与定价}

\subsection{Anthropic 生态的吸引力}

作者认为选择编程代理不仅是选工具，更是选生态：

\begin{itemize}
    \item \textbf{Anthropic 生态}：Claude Chat + Claude Code + Claude Cowork，形成类似 Apple 的闭环体验。Anthropic 还在逐步构建 OpenClaw（主动式个人代理）的安全版本。
    \item \textbf{OpenAI 生态}：除 Codex 外，其他产品对作者吸引力不大。ChatGPT 在 UI、对话风格和模型选择上已落后于 Claude Chat。
\end{itemize}

\subsection{定价方案对比}

\begin{table}[H]
\centering
\caption{定价方案对比}
\begin{tabular}{lcc}
\toprule
\textbf{层级} & \textbf{Claude Code} & \textbf{Codex} \\
\midrule
入门版 & \$20/月 & \$20/月 \\
中间版（Max 5x） & \$100/月 & --- \\
重度使用版 & \$200/月 & \$200/月 \\
\bottomrule
\end{tabular}
\end{table}

\begin{importantbox}{Claude Code 的定价优势}
Claude Code 提供了 \$100/月 的中间层级（Max 5x），对大多数开发者来说已经足够。这避免了从 \$20 直接跳到 \$200 的尴尬，使其\textbf{实际使用成本更低}。
\end{importantbox}

\subsection{Skills 与 Plugins 生态}

\begin{itemize}
    \item Skills 在两个平台间兼容，体验一致
    \item Codex 的 plugin 支持推出较晚，可用插件较少
    \item 社区内容（Reddit、X、博客）以 Claude Code 为主，反映其更大的社区规模
    \item 作者认为大多数开发者并不使用 plugins，这不应成为决策的关键因素
\end{itemize}

\subsection{本章小结}

Anthropic 正在构建一个紧密整合的产品生态，对于已经使用 Claude Chat 的用户吸引力尤其大。Claude Code 的三级定价比 Codex 的两级定价更灵活。Skills 生态基本通用，plugins 差异不大。

